\chapter{Electrónica, software e instrumentación}
\label{cap:electronica}

\section{Distribución de responsabilidades}
La electrónica diferencia el control de plataforma, la conectividad y la actuación del manipulador. El ATmega2560 atiende motores DC, encoders, línea, ultrasonido, barrido del radar, batería e interfaz local. El ESP8266 sirve la aplicación Web y traduce solicitudes entre WebSocket y UART. El control del brazo debe corresponder al bus TTL serie de los servos STS y a sus tensiones de alimentación.

Los servos de bus no se controlan como servos RC por pulsos individuales desde un PCA9685. La integración requiere definir transmisión semidúplex, referencia eléctrica, alimentación, límites y manejo de fallos. La presencia de un componente en el modelo CAD representa ocupación y montaje, sin demostrar su integración funcional.

\begin{table}[!htbp]\centering\small
\caption{Asignación principal de señales en el Mega.}
\begin{tabularx}{\linewidth}{@{}P{38mm}YP{37mm}@{}}
\toprule\textbf{Elemento} & \textbf{Señales} & \textbf{Función}\\\midrule
Motor delantero derecho & D22, D24; PWM D9 & Puente H.\\
Motor delantero izquierdo & D26, D28; PWM D10 & Puente H.\\
Motor trasero derecho & D5, D6; PWM D11 & Puente H.\\
Motor trasero izquierdo & D7, D8; PWM D12 & Puente H.\\
Hall FR / FL & A8/A12; A9/A13 & Canales A/B.\\
Hall RR / RL & A10/A14; A11/A15 & Canales A/B.\\
Línea IR1 a IR5 & A4, A3, A2, A1, A0 & Lectura digital configurada.\\
Batería & A5 & Divisor resistivo.\\
Ultrasonido & TRIG D30; ECHO D31 & Tiempo de vuelo.\\
Servo de radar & D13 & Consigna de barrido.\\
OLED & SDA D20; SCL D21 & Único dispositivo I2C.\\
Encoder de menú & CLK D2; DT D3; SW D42 & Giro y pulsación.\\
Enlace ESP & TX1 D18; RX1 D19 & UART de operación.\\
Servicio ESP & D32 KEY; D33 RESET & Accionamiento mediante NPN.\\\bottomrule
\end{tabularx}
\end{table}

Los motores reciben energía desde los puentes H. Los GPIO transportan señales de control y no sustituyen las etapas de potencia. La masa lógica es común; una salida negativa de motor del puente H no se usa como referencia de tierra. El orden físico de los conectores Hall debe convertirse al orden lógico de ruedas antes de aplicar la cinemática.

\section{Alimentación y medida de batería}
La distribución debe considerar los motores de la base, lógica, radio, radar y servos del brazo. Las referencias nominales usadas en la selección son 12 V para STS3250 y 6 V para STS3032. La potencia debe dimensionarse por simultaneidad, arranque, caída de tensión y disipación. Para ramales de tensión $V_j$, corriente $I_j$ y eficiencia de conversión $\eta_j$:
\begin{align}
P_{\mathrm{carga}}&=\sum_j V_j I_j,\\
I_{\mathrm{bat}}&\simeq\frac{1}{V_{\mathrm{bat}}}\sum_j\frac{V_jI_j}{\eta_j}.
\end{align}
La autonomía aproximada se expresa como
\begin{equation}
t_{\mathrm{aut}}\simeq\frac{E_{\mathrm{útil}}}{P_{\mathrm{media}}},
\qquad E_{\mathrm{nom}}=V_{\mathrm{nom}}C_{\mathrm{Ah}}.
\end{equation}
La energía se expresa en Wh cuando la capacidad está en Ah. El valor útil depende de corriente, tensión de corte, envejecimiento y conversión; no se asigna una autonomía sin medir esas condiciones.

La batería se adquiere por A5 mediante dos resistencias nominales de $10\unit{k\Omega}$ y un capacitor de $100\unit{nF}$. Para una lectura $a$ de un ADC de $n$ bits:
\begin{equation}
V_{\mathrm{bat}}=\frac{a}{2^n-1}V_{\mathrm{ref}}
\frac{R_{\mathrm{sup}}+R_{\mathrm{inf}}}{R_{\mathrm{inf}}}.
\label{eq:bateria}
\end{equation}
Con resistencias iguales, la escala nominal es dos. La impedancia de Thévenin y la constante de tiempo son
\begin{equation}
R_{\mathrm{th}}=\frac{R_{\mathrm{sup}}R_{\mathrm{inf}}}{R_{\mathrm{sup}}+R_{\mathrm{inf}}},
\qquad \tau_{RC}=R_{\mathrm{th}}C=0{,}5\unit{ms}.
\end{equation}
El capacitor filtra ruido, pero no protege de sobretensión. El error de referencia, la tolerancia del divisor y las caídas de masa afectan la lectura. La tensión no equivale a un porcentaje de carga universal.

\section{Modelo de los motores DC}
El modelo electromecánico equivalente relaciona tensión $v_m$, corriente $i$, velocidad $\omega_m$ y par de carga $\tau_L$. Con resistencia $R_m$, inductancia $L_m$, constante de fuerza contraelectromotriz $K_e$, constante de par $K_t$, inercia $J_m$ y rozamiento viscoso $b_m$:
\begin{align}
v_m&=R_mi+L_m\dot i+K_e\omega_m,\label{eq:motor1}\\
J_m\dot\omega_m&=K_ti-b_m\omega_m-\tau_L.\label{eq:motor2}
\end{align}
Si $\vct{x}_m=[i\;\omega_m]^{\trans}$, ambas ecuaciones forman el sistema
\begin{equation}
\dot{\vct{x}}_m=
\underbrace{\begin{bmatrix}
-R_m/L_m&-K_e/L_m\\
K_t/J_m&-b_m/J_m
\end{bmatrix}}_{\mat{A}_m}\vct{x}_m+
\begin{bmatrix}1/L_m&0\\0&-1/J_m\end{bmatrix}
\begin{bmatrix}v_m\\\tau_L\end{bmatrix}.
\end{equation}
Los parámetros del motor instalado no están identificados; el modelo explica la estructura de su respuesta. En un puente H ideal, el PWM produce una tensión media $\overline v_m\simeq sDV_{\mathrm{bus}}$, con ciclo útil $0\le D\le1$ y sentido $s=\pm1$. Las pérdidas, zona muerta, carga y batería impiden interpretar PWM como velocidad medida.

Para una reducción $i_r=\omega_m/\omega_w$, la velocidad de rueda es $\omega_w=\omega_m/i_r$ y el par de salida se aproxima por $\tau_w=\eta i_r\tau_m$. Una inercia de carga referida al motor se divide por $i_r^2$.

\section{Sensores de entorno}
\subsection{Ultrasonido y barrido}
La distancia acústica $d_s$ se calcula a partir del tiempo de ida y vuelta $\Delta t$ y velocidad del sonido $c$:
\begin{equation}
d_s=\frac{c\Delta t}{2},
\qquad c\simeq331{,}3+0{,}606\,T_C\quad[\mathrm{m/s}],
\end{equation}
donde $T_C$ está en grados Celsius. A $20\,^{\circ}\mathrm C$, $c\simeq343{,}4\unit{m/s}$. Una cuantización temporal de $25\unit{\mu s}$ produce un incremento nominal de aproximadamente $4{,}29\unit{mm}$; ese incremento no representa exactitud.

Para un ángulo de barrido $\alpha$ y origen del sensor conocido, la proyección ideal en su plano es
\begin{equation}
\vct{p}_s=\begin{bmatrix}d_s\cos\alpha\\d_s\sin\alpha\\0\end{bmatrix},
\qquad \vct{p}_B=\mat{R}_{Bs}\vct{p}_s+\vct{t}_{Bs}.
\end{equation}
La consigna de servo no es una medida de su ángulo. El eco depende de superficie, inclinación, haz y tiempo de establecimiento. La muestra debe conservar validez y antigüedad. En el funcionamiento de plataforma considerado, el HC-SR04 se usa para observación y telemetría; no se le atribuye una función física de parada que no esté implementada.

\subsection{Barra de línea}
Los cinco sensores se interpretan como señales digitales en la configuración de trabajo. Si $b_i\in\{0,1\}$ representa presencia de línea y $x_i$ la posición transversal, el centroide es
\begin{equation}
e_{\ell}=\frac{\sum_{i=1}^5 b_i x_i}{\sum_{i=1}^5b_i},
\qquad \sum_i b_i>0.
\end{equation}
La implementación usa posiciones relativas $2,1,0,-1,-2$ y corrección proporcional de rumbo. Cuando no hay sensores activos o la muestra es inválida, el seguimiento se detiene. No se asigna error cero a una línea perdida.

\section{Software, comunicaciones e interfaz}
La cadena de control remoto se organiza en solicitud Web, transporte por ESP, recepción UART y ejecución supervisada por ATmega. El protocolo de enlace es ROBO2/perfil 2 con extensión EX01. Los mensajes deben conservar versión, unidades, longitud, correlación, resultado y vigencia. Una aceptación de comando confirma su admisión, mientras el fin de una misión se observa por el estado de ejecución.

El control se atiende nominalmente cada 10 ms. Una demora superior a 100 ms durante movimiento provoca corte de salidas; los incumplimientos repetidos tienen tratamiento adicional. Render de OLED, adquisición y tráfico comparten recursos, de modo que la suficiencia de memoria y tiempos debe evaluarse para cada configuración de software y carga.

La interfaz local tiene una OLED de $128\times64$ píxeles y un encoder. El corto confirma o alterna; desde el dashboard abre el menú completo. El largo se ejecuta al soltar después de al menos 800 ms: abre controles rápidos desde el dashboard y vuelve al contexto anterior desde el resto. Si existen cambios pendientes, la elección Sí/No es obligatoria.

La Web distingue datos reales, ausencia de conexión y simulación. Las secuencias locales pueden seguir ejecutándose sin una WebSocket abierta si el modo lo permite; las consignas manuales conservan su propia vigencia. El cargador Windows separa diagnóstico, operación y mantenimiento. El enlace de operación y servicio utiliza D18/D19; KEY y RESET del ESP pueden accionarse mediante NPN en D32/D33, manteniendo RESET del Mega libre en el servicio por software. La compatibilidad eléctrica debe comprobarse en el shield real.
\FloatBarrier
